\documentclass[11pt,letterpaper]{report}

\usepackage[utf8]{inputenc}
\usepackage[T1]{fontenc}
\usepackage{amsmath,amssymb,amsthm}
\usepackage{graphicx}
\usepackage{geometry}
\usepackage{hyperref}
\usepackage{listings}
\usepackage{xcolor}
\usepackage{booktabs}
\usepackage{array}
\usepackage{longtable}
\usepackage{caption}
\usepackage{subcaption}
\usepackage{enumitem}
\usepackage{fancyhdr}
\usepackage{titlesec}
\usepackage{minted}
\usepackage{float}
\usepackage{siunitx}
\usepackage{bm}
\usepackage{gensymb}

\geometry{margin=1in}
\pagestyle{fancy}
\fancyhf{}
\fancyhead[L]{SAFB C Library Documentation}
\fancyhead[R]{\thepage}
\fancyfoot[C]{\thepage}

\hypersetup{
    colorlinks=true,
    linkcolor=blue,
    filecolor=magenta,
    urlcolor=cyan,
}

\lstset{
    basicstyle=\ttfamily\small,
    keywordstyle=\color{blue}\bfseries,
    commentstyle=\color{green!60!black},
    stringstyle=\color{red!70!black},
    frame=single,
    breaklines=true,
    tabsize=4,
    showstringspaces=false,
}

\newcommand{\qvec}{\mathbf{q}}
\newcommand{\rvec}{\mathbf{r}}
\newcommand{\hvec}{\mathbf{h}}
\newcommand{\Rmat}{\mathbf{R}}
\newcommand{\Gmat}{\mathbf{G}}
\newcommand{\Ginv}{\mathbf{G}^{-1}}
\newcommand{\bfs}{\boldsymbol{\psi}}
\newcommand{\bfsf}{\bfs_f}
\newcommand{\bfsr}{\bfs_r}
\newcommand{\cvec}{\mathbf{c}}
\newcommand{\avg}[1]{\langle #1 \rangle}
\newcommand{\dof}{\mathrm{dof}}

\title{\textbf{Symmetry-Adapted Fourier Basis (SAFB) C Library}\\[1em]
Documentation and Theory Reference}
\author{SAFB Migration Project}
\date{\today}

\begin{document}

\maketitle

\tableofcontents

\chapter{Introduction to SAFB}

\section{Overview}

The Symmetry-Adapted Fourier Basis (SAFB) C library generates scalar fields that strictly adhere to crystallographic space group symmetry constraints. The library is designed for initializing Self-Consistent Field Theory (SCFT) simulations of block copolymer self-assembly, where proper symmetry-adapted initial conditions are critical for convergence.

\section{Key Concepts}

\begin{itemize}
    \item \textbf{Space Group}: A mathematical group of symmetry operations (rotations and translations) that describe the symmetry of a crystal structure. There are 230 space groups in 3D and 17 wallpaper groups in 2D.

    \item \textbf{Star}: A set of symmetry-equivalent reciprocal lattice vectors (Miller indices h,k,l). A star contains all vectors generated by applying the point group operations to a single seed vector.

    \item \textbf{SAFBBasis}: A collection of valid Fourier modes (stars) for a given space group and lattice. Each mode has an associated amplitude and phase determined by symmetry constraints.

    \item \textbf{Field Generation}: The process of computing the real-space field $\psi(\mathbf{r})$ from symmetry-adapted Fourier coefficients using inverse FFT.

    \item \textbf{VTK Export}: Writing the generated field to VTK XML format for visualization in ParaView or similar tools.
\end{itemize}

\section{Use Cases}

\begin{enumerate}
    \item \textbf{SCFT Initialization}: Generate symmetry-adapted initial conditions for block copomer self-consistent field theory simulations.
    \item \textbf{Crystal Structure Analysis}: Study the symmetry properties of different crystal structures (gyroid, BCC, FCC, etc.).
    \item \textbf{Educational Tool}: Visualize and understand the relationship between space group symmetry and Fourier mode structure.
\end{enumerate}

\section{Supported Space Groups}

The library supports all 230 3D space groups and all 17 2D wallpaper groups. The 3D space groups are stored in \texttt{examples/space\_groups\_3d/} and the 2D wallpaper groups in \texttt{examples/space\_groups\_2d/}.

\subsection{Notable 3D Space Groups}

\begin{itemize}
    \item \textbf{Pm-3m} (No. 221): Simple cubic, used for BCC-like morphologies
    \item \textbf{Ia-3d} (No. 230): Body-centered cubic gyroid, the most studied bicontinuous structure
    \item \textbf{Pn-3m} (No. 224): Diamond structure
    \item \textbf{P4\_132} (No. 213): Chiral cubic structure
\end{itemize}

\subsection{Supported 2D Wallpaper Groups}

\begin{itemize}
    \item \textbf{Oblique}: P1, P2 (2 groups)
    \item \textbf{Rectangular}: Pm, Pg, Pmm, Pmg, P2gg (5 groups)
    \item \textbf{Square}: P4, P4mm, P4gm, P422 (4 groups)
    \item \textbf{Hexagonal}: P3, P3m1, P31m, P6, P6mm (5 groups)
\end{itemize}

\chapter{Crystallographic Background}

\section{Space Groups}

A space group is a mathematical description of the symmetry of a crystal structure. It consists of all symmetry operations (rotations, reflections, inversions, and translations) that leave the crystal invariant.

\subsection{Symmetry Operations}

Each symmetry operation consists of a rotation matrix $\Rmat$ and a fractional translation vector $\mathbf{t}$:

\begin{equation}
    \rvec' = \Rmat \cdot \rvec + \mathbf{t}
\end{equation}

where $\rvec$ is a position vector in direct space.

For 3D space groups, $\Rmat$ is a $3 \times 3$ integer matrix and $\mathbf{t}$ is a 3-vector of fractions. For 2D wallpaper groups, $\Rmat$ is a $2 \times 2$ integer matrix and $\mathbf{t}$ is a 2-vector of fractions (promoted to 3D with z-invariance).

\subsection{Point Group vs. Space Group}

\begin{itemize}
    \item The \textbf{point group} consists only of the rotation/reflection matrices (no translations). It describes the rotational symmetry.
    \item The \textbf{space group} includes both rotations and translations. It describes the full symmetry.
    \item The order of the point group (number of rotation matrices) is always a divisor of the order of the space group (number of symmetry operations).
\end{itemize}

\section{Reciprocal Lattice}

The reciprocal lattice is the Fourier transform of the direct lattice. For a direct lattice with basis vectors $\mathbf{a}_1, \mathbf{a}_2, \mathbf{a}_3$, the reciprocal lattice basis vectors are:

\begin{equation}
    \mathbf{b}_1 = 2\pi \frac{\mathbf{a}_2 \times \mathbf{a}_3}{\mathbf{a}_1 \cdot (\mathbf{a}_2 \times \mathbf{a}_3)}, \quad
    \mathbf{b}_2 = 2\pi \frac{\mathbf{a}_3 \times \mathbf{a}_1}{\mathbf{a}_1 \cdot (\mathbf{a}_2 \times \mathbf{a}_3)}, \quad
    \mathbf{b}_3 = 2\pi \frac{\mathbf{a}_1 \times \mathbf{a}_2}{\mathbf{a}_1 \cdot (\mathbf{a}_2 \times \mathbf{a}_3)}
\end{equation}

Any reciprocal lattice vector can be written as:

\begin{equation}
    \qvec = h \mathbf{b}_1 + k \mathbf{b}_2 + l \mathbf{b}_3 = (h, k, l)_{\text{recip}}
\end{equation}

where $(h, k, l)$ are the Miller indices (integers).

\subsection{Metric Tensor}

The metric tensor $\Gmat$ encodes the geometry of the reciprocal lattice:

\begin{equation}
    G_{ij} = \mathbf{b}_i \cdot \mathbf{b}_j
\end{equation}

The squared magnitude of a reciprocal lattice vector is:

\begin{equation}
    |\qvec|^2 = \mathbf{q}^T \Gmat \mathbf{q} = \sum_{i,j} h_i G_{ij} h_j
\end{equation}

\section{Stars}

A \textbf{star} of a reciprocal lattice vector $\qvec$ is the set of all vectors generated by applying the point group operations:

\begin{equation}
    \text{Star}(\qvec) = \{ \Rmat \cdot \qvec \mid \Rmat \in \text{Point Group} \} \cup \{ -\Rmat \cdot \qvec \mid \Rmat \in \text{Point Group} \}
\end{equation}

The union with negatives ensures the star is closed under inversion, which is required for real-valued Fourier series.

\subsection{Star Properties}

\begin{itemize}
    \item \textbf{Multiplicity}: The number of vectors in the star. Determined by the order of the point group divided by the order of the stabilizer subgroup.
    \item \textbf{Closure}: A star is closed if it contains all symmetry-equivalent vectors. The function \texttt{star\_is\_closed()} verifies this.
    \item \textbf{q-squared}: All vectors in a star have the same $|\qvec|^2$ (q-squared), since rotations preserve length.
\end{itemize}

\section{Miller Indices}

Miller indices $(h, k, l)$ label reciprocal lattice vectors and crystallographic planes. In the SAFB library, Miller indices are used to:

\begin{itemize}
    \item Identify and group reciprocal lattice vectors into stars
    \item Compute the canonical key for lexicographic sorting
    \item Label modes in the output
\end{itemize}

The canonical key is computed by sorting the Miller indices and taking the absolute value of the first non-zero element, ensuring a unique identifier for each family of planes.

\chapter{Mathematical Formulation}

\section{Symmetry-Adapted Fourier Series}

The scalar field $\psi(\rvec)$ is expanded as a Fourier series over symmetry-adapted modes:

\begin{equation}
    \psi(\rvec) = \sum_{n} A_n \exp(i \qvec_n \cdot \rvec)
\end{equation}

where:
\begin{itemize}
    \item $\qvec_n$ are the reciprocal lattice vectors (from stars)
    \item $A_n$ are complex amplitudes constrained by symmetry
\end{itemize}

\section{Phase Constraints}

Symmetry imposes constraints on the Fourier coefficients. For a symmetry operation $(\Rmat, \mathbf{t})$ and a reciprocal vector $\qvec$:

\begin{equation}
    c_{\Rmat \qvec} = \exp(2\pi i \, \mathbf{t} \cdot \qvec) \, c_{\qvec}
\end{equation}

This means that once the coefficient for one vector in a star is known, all other coefficients in the star are determined up to a global phase factor.

\subsection{Phase Factor Computation}

The phase factor for a symmetry operation $(\Rmat, \mathbf{t})$ acting on vector $\qvec = (h, k, l)$ is:

\begin{equation}
    \phi = 2\pi \, \mathbf{t} \cdot \qvec = 2\pi \sum_{j} t_j q_j
\end{equation}

where the dot product is computed using fractional coordinates. The exponential $\exp(i\phi)$ gives the phase relationship between symmetry-equivalent coefficients.

\section{Coefficient Solver}

The phase constraints form a system of linear equations. For a star with $m$ vectors and $r$ relationships, we have:

\begin{equation}
    c_j = e^{i\phi_j} c_k
\end{equation}

for each relationship $(j, k, \phi_j)$. The system is solved using BFS (Breadth-First Search) on the relationship graph:

\begin{enumerate}
    \item Start with a reference coefficient $c_{\text{ref}} = 1.0 + 0i$
    \item For each relationship, propagate the phase constraint through the graph
    \item Check for consistency (no contradictory phase constraints)
    \item All coefficients are determined up to a global phase factor
\end{enumerate}

\section{Normalization}

The field is normalized using the hyperbolic tangent function to ensure values lie in $[0, 1]$:

\begin{equation}
    \psi_{\text{norm}}(\rvec) = \frac{1}{2} \left[ 1 + \tanh\left( \frac{\psi(\rvec) - \avg{\psi}}{\sigma \cdot \text{std}(\psi)} \right) \right]
\end{equation}

where $\sigma$ is a scaling parameter (typically $\sigma = 1.0$).

\section{Square Norm}

The square norm (average of $\psi^2$ over the unit cell) is:

\begin{equation}
    \avg{|\psi|^2} = \frac{1}{V} \int_V |\psi(\rvec)|^2 \, d\rvec
\end{equation}

This can be computed numerically on the FFT grid or analytically using Parseval's theorem:

\begin{equation}
    \avg{|\psi|^2} = \sum_n |A_n|^2
\end{equation}

\chapter{Numerical Methods}

\section{FFT-Based Field Generation}

The field is computed using the inverse Fast Fourier Transform (FFT):

\begin{equation}
    \psi(\rvec_j) = \mathcal{F}^{-1}\{ c_{\qvec_n} \}
\end{equation}

\subsection{FFTW Integration}

The library uses FFTW (Fastest Fourier Transform in the West) for efficient FFT computation:

\begin{itemize}
    \item \texttt{fftw\_plan\_dft\_3d}: Creates a 3D FFT plan for the specified grid size
    \item \texttt{fftw\_execute}: Executes the FFT plan
    \item \texttt{fftw\_destroy\_plan}: Cleans up the plan
\end{itemize}

\subsection{Grid Quadrature}

The real-space grid is defined by the lattice parameters and the number of grid points:

\begin{equation}
    x_i = \frac{i}{N_x} L_x, \quad y_j = \frac{j}{N_y} L_y, \quad z_k = \frac{k}{N_z} L_z
\end{equation}

where $N_x, N_y, N_z$ are the grid dimensions and $L_x, L_y, L_z$ are the lattice lengths.

\section{Coefficient Solver Algorithm}

The coefficient solver uses BFS on the star's relationship graph:

\begin{algorithm}
\caption{BFS Coefficient Solver}
\begin{algorithmic}[1]
\State $c_{\text{ref}} \gets 1.0 + 0i$
\State $Q \gets [\text{ref\_index}]$
\State $\text{visited}[\text{ref\_index}] \gets \text{true}$
\While{$Q$ is not empty}
    \State $u \gets Q.\text{pop}()$
    \For{each neighbor $(v, \phi)$ of $u$}
        \If{$\text{visited}[v]$}
            \State Verify: $c_v == e^{i\phi} c_u$
        \Else
            \State $c_v \gets e^{i\phi} c_u$
            \State $\text{visited}[v] \gets \text{true}$
            \State $Q.\text{push}(v)$
        \EndIf
    \EndFor
\EndWhile
\end{algorithmic}
\end{algorithm}

\section{Memory Layout}

The library uses flat arrays for efficiency:

\begin{itemize}
    \item Reciprocal space: flat array of complex coefficients, indexed by star and vector position
    \item Real space: flat array of real field values, indexed by $(i, j, k)$ in row-major order
    \item The mapping between reciprocal and real space is handled by FFTW
\end{itemize}

\chapter{Library Architecture}

\section{Module Overview}

\begin{table}[h]
\centering
\begin{tabular}{lll}
\hline\hline
\textbf{Module} & \textbf{Purpose} & \textbf{Key Functions} \\
\hline
\texttt{domain.h/c} & Data structures & \texttt{LatticeInfo}, \texttt{Star}, \texttt{SAFBBasis} \\
\texttt{symmetry\_ops.h/c} & Symmetry operations & \texttt{read\_spacegroup\_ops\_txt}, \texttt{star\_from\_hkl} \\
\texttt{basis.h/c} & Basis construction & \texttt{family\_planes\_info}, \texttt{build\_basis} \\
\texttt{initializers.h/c} & Amplitude assignment & \texttt{random\_initialization}, \texttt{manual\_initialization} \\
\texttt{field.h/c} & Field generation & \texttt{build\_field}, \texttt{build\_field\_2d} \\
\texttt{engine.h/c} & High-level API & \texttt{engine\_create}, \texttt{engine\_init\_random} \\
\texttt{analytic.h/c} & Analytical calculations & \texttt{calculate\_square\_norm} \\
\hline\hline
\end{tabular}
\caption{Library modules and their responsibilities}
\end{table}

\section{Data Structures}

\subsection{LatticeInfo}

\begin{lstlisting}
typedef struct {
    double Lx, Ly, Lz;      // Lattice lengths
    double alpha, beta, gamma; // Lattice angles (degrees)
    double G[3][3];         // Reciprocal metric tensor
    double Ginv[3][3];      // Inverse metric tensor
    int dim;                // 2 or 3
} LatticeInfo;
\end{lstlisting}

\subsection{Star}

\begin{lstlisting}
typedef struct {
    int hkl[3];             // Canonical Miller indices
    double q2;              // |q|^2 (q-squared)
    int multiplicity;       // Number of vectors
    int closed;             // 1 if closed under inversion
    int vectors[3][MAX_STAR_SIZE]; // Symmetry-equivalent vectors
    int vector_keys[MAX_STAR_SIZE]; // Lexicographic keys
    double phase_factors[MAX_STAR_SIZE]; // Phase factors
} Star;
\end{lstlisting}

\subsection{SAFBBasis}

\begin{lstlisting}
typedef struct {
    int n_modes;            // Number of modes (families)
    Star modes[MAX_MODES];  // Mode data
    int centrosymmetric;    // 1 if inversion is in point group
    int inversion_at_origin; // 1 if inversion at origin
} SAFBBasis;
\end{lstlisting}

\section{API Usage}

\begin{lstlisting}
// Create engine from space group file
Engine *engine = engine_create("examples/space_groups_3d/ia_3d.txt",
                                "examples/lattice_cubic.txt");

// Random initialization
engine_init_random(engine, 42);

// Build field
engine_build_field(engine, 64, 64, 64);

// Output field to VTK
engine_output_field(engine, "gyroid.vts");

// Clean up
engine_destroy(engine);
\end{lstlisting}

\chapter{User Manual}

\section{Building the Library}

\begin{verbatim}
source /sandbox/setup_build.sh   # Activate build environment
make                             # Build all test programs
make test                        # Run all tests
\end{verbatim}

\section{Running Examples}

\subsection{3D Example: Ia-3d Gyroid}

\begin{verbatim}
# Build the library first
source /sandbox/setup_build.sh
make

# Run the gyroid example (uses space group file)
./examples/run_gyroid.sh
\end{verbatim}

\subsection{2D Example: P4 Square Lattice}

\begin{verbatim}
# Run the 2D P4 example
./examples/run_2d_p4.sh
\end{verbatim}

\section{Configuration}

\subsection{Lattice Parameters}

Lattice parameters are specified in \texttt{examples/lattice\_cubic.txt}:

\begin{verbatim}
a=10, b=10, c=10, alpha=90, beta=90, gamma=90
\end{verbatim}

For non-cubic lattices, modify the values accordingly.

\subsection{Grid Resolution}

The grid resolution is specified in the engine call:

\begin{verbatim}
engine_build_field(engine, Nx, Ny, Nz)
\end{verbatim}

Typical values: $64 \times 64 \times 64$ for production, $16 \times 16 \times 16$ for testing.

\section{Output Files}

\begin{itemize}
    \item \texttt{*.vts}: VTK XML structured point data files (viewable in ParaView)
    \item \texttt{*.dat}: Field data in plain text format
    \item \texttt{*.json}: Basis metadata (mode count, q-values, etc.)
\end{itemize}

\chapter{Examples and Workflows}

\section{Gyroid (Ia-3d)}

The gyroid is a bicontinuous minimal surface with Ia-3d symmetry. It is one of the most studied structures in block copomer science.

\subsection{Generation Steps}

\begin{enumerate}
    \item Load space group operations from \texttt{examples/space\_groups\_3d/ia\_3d.txt}
    \item Load lattice parameters (cubic, $a = 10$)
    \item Build basis: find all valid stars up to $N_{\text{max}} = 5$
    \item Random initialization: assign random amplitudes with symmetry constraints
    \item Build field: compute inverse FFT on $64^3$ grid
    \item Normalize: apply tanh normalization
    \item Export: write to \texttt{gyroid.vts}
\end{enumerate}

\subsection{Expected Results}

\begin{itemize}
    \item 5 modes (families) with multiplicities: 8, 6, 24, 12, 48
    \item Total Fourier modes: 98
    \item Field range: $[0, 1]$ after normalization
    \item Square norm: $\approx 0.5$ for random initialization
\end{itemize}

\section{BCC-like (Pm-3m)}

The Pm-3m space group generates BCC-like density modulations.

\subsection{Generation Steps}

Same as gyroid, but with \texttt{pm3m.txt} space group file and cubic lattice.

\section{2D Examples}

\subsection{Square Lattice (P4)}

The P4 wallpaper group generates 4-fold symmetric patterns.

\begin{itemize}
    \item 4 symmetry operations (0°, 90°, 180°, 270° rotations)
    \item Supports 2D field generation
    \item Useful for studying 2D self-assembly patterns
\end{itemize}

\subsection{Hexagonal Lattice (P6)}

The P6 wallpaper group generates 6-fold symmetric patterns.

\begin{itemize}
    \item 6 symmetry operations (0°, 60°, 120°, 180°, 240°, 300° rotations)
    \item Supports 2D field generation
    \item Hexagonal packing patterns
\end{itemize}

\chapter{Glossary}

\section*{Glossary of Crystallographic Terms}

\begin{description}
\item[Space Group] A mathematical group describing all symmetry operations (rotations, reflections, inversions, and translations) that leave a crystal structure invariant. There are 230 space groups in 3D and 17 wallpaper groups in 2D.

\item[Wallpaper Group] A 2D space group describing the symmetry of patterns that repeat in two dimensions. Also called plane groups. There are 17 wallpaper groups.

\item[Point Group] The set of rotational and reflectional symmetry operations of a crystal, ignoring translations. The point group describes the rotational symmetry of the crystal class.

\item[Symmetry Operation] A transformation (rotation, reflection, inversion, or translation) that maps the crystal onto itself. Each operation is represented by a rotation matrix $\Rmat$ and a fractional translation vector $\boldsymbol{t}$.

\item[Miller Indices] Integer triplet $(h, k, l)$ that labels a reciprocal lattice vector or a family of crystallographic planes. The indices are the reciprocals of the intercepts of the plane with the crystallographic axes.

\item[Star] A set of symmetry-equivalent reciprocal lattice vectors. A star contains all vectors generated by applying the point group operations to a seed vector, plus their negatives.

\item[Multiplicity] The number of vectors in a star. Determined by the order of the point group divided by the order of the stabilizer subgroup of the seed vector.

\item[Reciprocal Lattice] The Fourier transform of the direct (real-space) lattice. Reciprocal lattice vectors $\qvec = h\boldsymbol{b}_1 + k\boldsymbol{b}_2 + l\boldsymbol{b}_3$ label the Fourier modes in the crystal.

\item[Direct Lattice] The real-space lattice defined by basis vectors $\boldsymbol{a}_1, \boldsymbol{a}_2, \boldsymbol{a}_3$. The crystal structure is a periodic arrangement of atoms on this lattice.

\item[Metric Tensor] A symmetric matrix $\Gmat$ encoding the geometry of the reciprocal lattice: $G_{ij} = \boldsymbol{b}_i \cdot \boldsymbol{b}_j$. Used to compute $|\qvec|^2$ for any Miller index triplet.

\item[Phase Constraint] A relationship between Fourier coefficients of symmetry-equivalent vectors: $c_{\Rmat \qvec} = \exp(2\pi i \, \boldsymbol{t} \cdot \qvec) \, c_{\qvec}$. Ensures the field respects space group symmetry.

\item[Fourier Mode] A single term in the Fourier expansion of the field: $A_n \exp(i \qvec_n \cdot \rvec)$. Each mode corresponds to a reciprocal lattice vector.

\item[SAFBBasis] A collection of valid Fourier modes (stars) for a given space group and lattice. Contains all the information needed to generate a symmetry-adapted field.

\item[Field Generation] The process of computing the real-space field $\psi(\rvec)$ from Fourier coefficients using inverse FFT.

\item[Normalization] The process of scaling the field to a standard range, typically $[0, 1]$, using a tanh function.

\item[Square Norm] The average of $\psi^2$ over the unit cell: $\avg{|\psi|^2} = \frac{1}{V} \int_V |\psi|^2 d\rvec$. Measures the overall amplitude of the field.

\item[Centrosymmetric] A crystal structure is centrosymmetric if it contains an inversion center. In terms of the point group, this means the inversion matrix $(-I)$ is a member.

\item[Inversion at Origin] A special case where the inversion operation maps $\rvec \to -\rvec$ (no translation). This simplifies the phase constraints.

\item[Canonical Key] A lexicographic identifier for a family of planes, computed by sorting the Miller indices and taking the absolute value of the first non-zero element.

\item[Closed Star] A star that contains all symmetry-equivalent vectors, including all negatives. Required for real-valued Fourier series.

\item[BFS Coefficient Solver] A breadth-first search algorithm that propagates phase constraints through the star's relationship graph to determine all Fourier coefficients.

\item[FFTW] Fastest Fourier Transform in the West. A widely-used library for computing FFTs. The SAFB library uses FFTW for efficient field generation.

\item[VTK] Visualization Toolkit. A file format for scientific visualization data. The SAFB library exports fields in VTK XML structured point data format.

\item[SCFT] Self-Consistent Field Theory. A mean-field theory for polymer self-assembly. The SAFB library provides symmetry-adapted initial conditions for SCFT simulations.

\item[Block Copolymer] A polymer consisting of two or more different polymer blocks covalently bonded together. Block copolymers self-assemble into periodic nanostructures.

\item[Gyroid] A bicontinuous minimal surface with Ia-3d symmetry. One of the most studied structures in block copolymer science.

\item[BCC] Body-centered cubic. A crystal structure with space group Pm-3m. Block copolymers can form BCC-like density modulations.

\item[FCC] Face-centered cubic. A crystal structure with space group Fm-3m.

\item[Diamond] A crystal structure with space group Pn-3m. Block copolymers can form diamond-like structures.

\item[Wallpaper Group P1] The simplest 2D space group, containing only the identity operation. No symmetry beyond translation.

\item[Wallpaper Group P6] A 2D space group with 6-fold rotational symmetry. Generates hexagonal patterns.

\item[Wallpaper Group P4] A 2D space group with 4-fold rotational symmetry. Generates square patterns.

\end{description}


\chapter{References}

\bibliographystyle{plain}
\bibliography{references}

\end{document}
